\section{From local occupations to the constant mode}\label{sec:assembly}

The energy calculation determines the occupation outside cellwise constants
at scale $R$. We now show that the remaining difference between cellwise
constants and the global constant is negligible on the depletion scale.
The distant comparison controls variation of cell root mean squares;
the local fourth-moment comparison lets us replace those root mean squares
by the arithmetic means that define orthogonal projections.

\subsection{A variance bound from distant pairs}

Let $s_v$ be the nonnegative bath function defined by
\eqref{eq:palm-density}, for $v\in\cC_R$. We regard these functions as a
finite array with values in the Hilbert space $L^2$ of the $N-1$ bath
coordinates. Its mean is $\bar s=\E_{v\in\cC_R}s_v$.

\begin{lemma}[Distant-pair spectral gap]\label{lem:distant-gap}
Let $M$ tend to infinity through integers, and let $(f_z)_{z\in(\Z/M\Z)^3}$
be an array in a Hilbert space. Average pairs $(z,z+h)$ uniformly over
$z$, over all transverse displacements, and over first-coordinate
displacements $h_1$ in $[.30M,.45M]$. For all sufficiently large $M$,
\[
 \E_z\|f_z-\bar f\|^2\le C\E_{z,h}\|f_{z+h}-f_z\|^2
\]
with an absolute constant $C$.
\end{lemma}
\begin{proof}
Use the finite Fourier transform of the array and Parseval's identity.
The multiplier of the right side is the displacement average of
$|e^{2\pi i k\cdot h/M}-1|^2$. If either transverse component of $k$
is nonzero, the average phase vanishes and this multiplier equals $2$.
For $k=(k_1,0,0)$, geometric summation gives
\[
 \left|\E_{h_1}e^{2\pi i k_1h_1/M}\right|
 \le \frac{C}{\min(k_1,M-k_1)},\qquad 1\le k_1<M.
\]
This gives a positive lower bound when the denominator exceeds a fixed
constant. For each of the finitely many remaining nonzero frequencies
(measured from either end of the frequency interval), the multiplier
converges to the average of $|e^{2\pi i kt}-1|^2$ over
$t\in[.30,.45]$. That integral is strictly positive. The finite minimum
is positive, completing the proof also for Hilbert-valued arrays.
\end{proof}

Apply the lemma with $M=K/R$ and use~\eqref{eq:distant-comparison}.
The result is
\begin{equation}\label{eq:rms-variance}
 \E_{v\in\cC_R}\|s_v-\bar s\|^2=o(D^{-1})+o_K(1).
\end{equation}
The scale-independent spectral gap here is the reason to average
macroscopically separated pairs. Nearest-neighbor squared differences
alone would introduce a volume-dependent factor.

\subsection{Arithmetic means and root mean squares}

Put $F(y,Y')=\sqrt V\,\Phi(y,Y')$, where $Y'$ denotes the bath
coordinates. The first-coordinate spatial measure is now normalized as
$\dd y/V$. For every cell $v$ define
\[
 m_v(Y')=\frac1{|v|}\int_v F(y,Y')\dd y,\qquad
 s_v(Y')=\left(\frac1{|v|}\int_v F(y,Y')^2\dd y\right)^{1/2}.
\]
The second definition agrees with~\eqref{eq:palm-density}. The quantities
$m_v$ are precisely the cellwise values used by $P_l$. Positivity of
$\Phi$ gives $0\le m_v\le s_v$ and
$(s_v-m_v)^2\le s_v^2-m_v^2$. For microcells, integrating this inequality
and using Proposition~\ref{prop:local-energy} gives
\begin{equation}\label{eq:micro-rms}
 \E_{v\in\cC_d}\|s_v-m_v\|^2
 \le\frac1N\langle\dGamma(W_d)\rangle_\Phi
 \le D^{-1}(\epsilon_j+o(1)),\qquad \epsilon_j\longrightarrow0.
\end{equation}
We must transfer this estimate to an $R$-cell. Averaging the $m_v$ over
its microcells gives its arithmetic mean exactly; averaging the $s_v$
does not give its root mean square. The next inequality controls that
discrepancy with the fourth-order information supplied by
\eqref{eq:local-fourth}.

\begin{lemma}[Root mean square comparison]\label{lem:rms-comparison}
For a finite nonnegative array $(y_i)$, with uniform independent indices
$i,i'$, one has
\begin{equation}\label{eq:rms-fourth-inequality}
 \left(\sqrt{\E_i y_i^2}-\E_i y_i\right)^2
 \le\frac12\E_{i,i'}
          \frac{|y_i-y_{i'}|^4}{y_i^2+y_{i'}^2},
\end{equation}
where the quotient is zero when its denominator vanishes. Moreover,
\begin{equation}\label{eq:square-root-quotient}
 |\sqrt y-\sqrt z|^4\le\frac{|y-z|^4}{y^2+z^2}
 \le8|\sqrt y-\sqrt z|^4\qquad(y,z\ge0).
\end{equation}
\end{lemma}
\begin{proof}
Write $A=\E y_i^2$ and $b=\E y_i$. If $A=0$, all assertions are
immediate. Otherwise Cauchy--Schwarz gives
\[
 \frac12\E_{i,i'}\frac{|y_i-y_{i'}|^4}{y_i^2+y_{i'}^2}
 \ge\frac{(A-b^2)^2}{A}
 \ge(\sqrt A-b)^2.
\]
For the second assertion, factor $y-z=(\sqrt y-\sqrt z)(\sqrt y+\sqrt z)$.
The ratio $(\sqrt y+\sqrt z)^4/(y^2+z^2)$ lies between $1$ and $8$;
the zero cases follow directly.
\end{proof}

Within an $R$-cell $b$, apply this lemma pointwise in $Y'$ to the
microcell values $s_v$. Their root mean square is exactly $s_b$.
Join any two microcells of $b$ by at most $C_jR$ adjacent steps inside
$b$, telescope their square roots, and use
$(\sum_{k=1}^m |a_k|)^4\le m^3\sum_{k=1}^m|a_k|^4$.
Even bounding each path sum by the sum over every internal edge gives,
after averaging $b$ and using~\eqref{eq:local-fourth},
\begin{equation}\label{eq:coarse-rms-comparison}
 \E_{b\in\cC_R}
 \left\|s_b-\E_{v\subset b,\ v\in\cC_d}s_v\right\|^2
 \le C_jR^6D^{-1-u/8}=o(D^{-1}).
\end{equation}
Here $R^3$ comes from the fourth-power path inequality and another $R^3$
from the number of internal edges. The prescribed $w<u/10000$ ensures
$6w<u/8$. Jensen's inequality and~\eqref{eq:micro-rms} now give
\begin{equation}\label{eq:coarse-mean-rms}
 \E_{b\in\cC_R}\|s_b-m_b\|^2
 \le D^{-1}(2\epsilon_j+o(1)).
\end{equation}

\subsection{Completion of the depletion asymptotic}

Let $P_0$ be the projection onto the global constant function on the
scaled torus, and let a subscript $1$ indicate action in the first
particle coordinate. Centering a Hilbert-valued array is an orthogonal
projection and hence a contraction. Combining
\eqref{eq:rms-variance} and~\eqref{eq:coarse-mean-rms}, we obtain
\begin{equation}\label{eq:low-mode}
 \begin{split}
 D\|(P_R-P_0)_1\Phi\|^2
 &=D\E_{b\in\cC_R}\|m_b-\bar m\|^2\\
 &\le4\epsilon_j+o(1)+o_K(1).
 \end{split}
\end{equation}
Multiplication of an inner-limit remainder by $D$ is harmless: at fixed
$\rho$, $D$ converges to the fixed number $D_*$. Since $j$ can be
chosen arbitrarily large before the dilution threshold is taken,
\eqref{eq:low-mode} makes the very-low-frequency contribution negligible.

The projections $W_R$ and $P_R-P_0$ are orthogonal and sum to $1-P_0$.
Permutation symmetry and Proposition~\ref{prop:local-energy} therefore
yield, in the stipulated order of limits,
\begin{equation}\label{eq:scaled-final-depletion}
 \begin{split}
 D\bigl(1-\|(P_0)_1\Phi\|^2\bigr)
 &=\frac1V\langle\dGamma(W_R)\rangle_\Phi
       +D\|(P_R-P_0)_1\Phi\|^2\\
 &=\frac8{3\sqrt\pi}\alpha^{3/2}+o(1).
 \end{split}
\end{equation}
The equality of the leading term includes both bounds: the second
summand is nonnegative, tends to zero, and the first already has the
asymptotic displayed in Proposition~\ref{prop:local-energy}.

\begin{proof}[Proof of Theorem~\ref{thm:main}]
The coordinate rescaling, with wave-function factor $\sigma^{3N/2}$,
is unitary and maps constant functions to constant functions. It preserves
the occupation fraction, so
$B_\Gamma=\|(P_0)_1\Phi\|^2$ for the unique ground-state density matrix.
In the inner thermodynamic limit,
\[
 D\sqrt{\rho a_v^3}\longrightarrow
 D_*\sqrt{\rho a_v^3}=\alpha_0^{3/2},\qquad
 \alpha\longrightarrow\alpha_0>0.
\]
Divide~\eqref{eq:scaled-final-depletion} by the first quantity. For a
given accuracy, first choose $j$ to make the contribution from
$\epsilon_j$ sufficiently small, and then take $D_*$ sufficiently large
for all dilute errors and fixed-parameter thresholds. By
\eqref{eq:physical-base-scale}, this amounts to choosing
$\rho<\rho_0(v,\epsilon)$. The estimates hold on a tail of every
thermodynamic sequence, with no requirement on its rate of convergence.
They prove precisely the limsup assertion in Theorem~\ref{thm:main}.
\end{proof}
